\section{A Gaussian experiment with a genuine temporal penalty}
\label{sec:v9-delayed}
The append-register Gaussian theorem in the supporting development
concerns a terminal vector experiment. The following experiment instead
forces information to survive until the question that uses it arrives.
Its final target is scalar, but its causal exponent depends on the
number of coordinates acquired before that question.

Fix $d\ge2$ and $\sigma>0$. Let $\theta_1,\ldots,\theta_d$ be independent
and uniform on $[-1,1]$, and observe sequentially
\[
                  X_j=\theta_j+\sigma Z_j,\qquad 1\le j\le d,
\]
where the $Z_j$ are independent standard Gaussians. After these
observations have been discarded, an independent uniform query
$J\in\{1,\ldots,d\}$ is revealed. The machine then updates its register
and estimates $\theta_J$. There is no external observation transcript.
The query is an input to that final update, not an additional persistent
decoder register. The cost is squared loss, bounded by four after
projecting the answer onto $[-1,1]$. Write
\[
 m(x)=\E[\theta_1\mid X_1=x],\qquad
 \beta=\E\Var(\theta_1\mid X_1).
\]
Let $\mathcal E_S$ be the infimal excess above $\beta$ when the total
register has at most $S$ states at every stage. Independent shared
randomness is permitted; the same conclusions hold without it.

\begin{theorem}[Delayed Gaussian query law]\label{thm:v9-delayed}
There are constants $0<c_{d,\sigma}\le C_{d,\sigma}<\infty$ such that
for every $S\ge2^d$,
\begin{equation}
             c_{d,\sigma}S^{-2/d}\le\mathcal E_S
                          \le C_{d,\sigma}S^{-2/d}.
 \label{eq:v9-delayedrate}
\end{equation}
The same order holds for the maximum over the $d$ individual coordinate
query risks, with a common encoder before the query.

In contrast, an encoder which receives the query together with the full
observation vector at its sole checkpoint has excess of order $S^{-2}$.
Thus the scalar terminal target has a strictly worse cardinality
exponent under the delayed-query causal constraint. The lower bound
uses the register \emph{before} the query and is not a final-output
dimension bound.
\end{theorem}
\begin{proof}
The posterior density on $[-1,1]$ is proportional to
$\exp(-(x-u)^2/(2\sigma^2))$. Differentiation under the integral on this
compact interval gives
\begin{equation}
                     m'(x)=\frac{\Var(\theta_1\mid X_1=x)}{\sigma^2}>0.
 \label{eq:v9-posteriorderivative}
\end{equation}
The posterior variance is positive because its density is positive on
an interval. The function $m$ is smooth, odd, and strictly increasing.
Choose $b>0$ and put $a=m(b)>0$. The density of $X_1$ is positive and
continuous, and $m'$ is positive and continuous on $[-b,b]$. The
change-of-variables formula therefore gives a density for $m(X_1)$
which is bounded below by a positive constant on $[-a,a]$. Independence
implies that $M=(m(X_1),\ldots,m(X_d))$ has density bounded below by a
positive constant on the cube $Q=[-a,a]^d$.

We record an elementary volume bound. If $c_1,\ldots,c_S$ are arbitrary
points in $\R^d$, and $v_d$ is the volume of the unit ball, choose
$r=(|Q|/(2Sv_d))^{1/d}$. The union of their radius-$r$ balls has volume
at most $|Q|/2$. Consequently
\begin{equation}
 \int_Q\min_{i\le S}\|x-c_i\|^2\,dx
       \ge\frac{|Q|}{2}\left(\frac{|Q|}{2Sv_d}\right)^{2/d}.
 \label{eq:v9-volumebound}
\end{equation}
This bound does not require cells to be convex or centres to lie in $Q$.

Consider any admissible machine and condition on its independent shared
random seed, when present. Its pre-query register $I_d$ has at most
$S$ values. Subsequent private updates and readouts can do no better than
being given $(I_d,J)$ and using its conditional posterior mean. Let
$c_i=(\E[\theta_j\mid I_d=i])_{j=1}^d$, with the seed included in all
these conditional expectations. The Markov relation from the full
observation to the register, and independence of the random seed,
imply $c_i=\E[M\mid I_d=i]$. Orthogonality of conditional expectation
shows that the resulting excess for the uniform query is
\[
              \frac1d\E\|M-c_{I_d}\|^2
       \ge \frac1d\E\min_{i\le S}\|M-c_i\|^2.
\]
For each fixed seed the law of $M$ is unchanged. Its density lower bound
on $Q$ and \eqref{eq:v9-volumebound} give
$c_{d,\sigma}S^{-2/d}$. Averaging the seed preserves this lower bound.
The argument permits randomized private encoding: replacing its selected
centre by the nearest centre is only a lower bound, not an assumption
that the encoding is deterministic. The maximum coordinate risk is at
least the uniform average, proving its lower bound as well.

For the upper bound put $L=\lfloor S^{1/d}\rfloor\ge2$ and partition
$[-1,1]$ into $L$ equal intervals. At observation $j$ compute the interval
label of $m(X_j)$ and append it in base $L$ to the stored labels. At that
time there are at most $L^j\le S$ states. After the query, replace the
whole label by the interval label of the requested coordinate; the
terminal decoder outputs its midpoint. This last update leaves only
$L$ possible labels and retains no separate query. Its squared excess
in every coordinate is at most $L^{-2}\le4S^{-2/d}$. Conditional
orthogonality again subtracts the same full-observation baseline $\beta$.
The construction is deterministic and works with or without a shared
seed.

For the noncausal checkpoint comparison, the encoder sees $(X,J)$ and
quantizes only $m(X_J)$ into $S$ equal intervals. This has excess at most
$S^{-2}$. The one-dimensional version of the same density lower bound
and \eqref{eq:v9-volumebound} gives the matching lower bound, even with
randomized encoding and a shared seed. Its $S$ terminal centres are
arbitrary real answers in $[-1,1]$, so no additional discreteness
assumption has entered. This proves all assertions.
\end{proof}

The theorem is a Bayesian slice of the parameterized Gaussian experiment
with the specified product prior, not a claim about uniform deficiency
over the whole Gaussian family. Its constants depend on $d$ and
$\sigma$. It controls persistent cardinality and finite acquisition
count, not the arithmetic cost of evaluating the real posterior mean.
Under a finite numerical signature, an approximation of $m$ requires a
separate accuracy and computation bound. The obstruction nevertheless
holds for every such implementation, since it was proved for the larger
unrestricted-arithmetic class.

The lower bound is the same conditional-variance geometry as
Corollary~\ref{cor:v9-quadratic}, now evaluated directly for a continuous
history law. In the finite theorem, successor compatibility determines
which cells can be executed. Here the late query forces all coordinate
readouts to factor through the same pre-query cell, and the volume bound
computes its cost. This explains both the lossless append-only regime
and a regime where forgetting before a future question changes the
exponent, without identifying either one with static terminal dimension.
